\section{A coefficient-sensitive obstruction with minimal subsequent repairs}
\label{sec:136-coefficient-history-obstruction}

\paragraph{Scope and inputs.}
The archive locators are cote136, physical pp.~117,119,121,
footers~116,118,120, scans~117,119,121.  We use the checked
first-skeleton replacement, its free envelope $L(F)$, the
relative coefficient comparison, the low-cohomology calculation
for a free degree-four pair in TWO-TRUNCATED-LINKING-REPAIR,
and AQ-136-RELATIVE-H1.  All coefficient and algebra
calculations below are included.  The history starts with
\emph{two specified scaled even disks}.  Their production
from $L(F)$ by a minimal source construction is not asserted.
The subsequent degree-three, degree-four and degree-two
repairs do use the actual minimal integer order, namely two.
This is a prescribed-map preservation counterexample, not
a counterexample to original connectedness.

\paragraph{ARCH-136-COEFFICIENT-HISTORY-OBSTRUCTION.}
There is an actual two-truncated source with literal supplied
first skeleton, graded-free form groups, rational vertical
coefficient comparison, $H^0=S$, $H^1=0$, and integer-torsion
$H^2$, with maps to two copies of one admissible theory,
such that every old degree-three torsion pair has common-graph
target primitives for every integer modulus.  It is obtained
from two scaled even disks by a minimal degree-three pair
and then a minimal degree-four pair.  A subsequent minimal
degree-two pair has target extensions, but every extension
has a nonzero degree-three graph residue in unbounded weights.
The obstruction persists after the supplied RH1 iteration
retaining the maps and their first skeleton.  In this example
RH1 is genuinely needed: a new one-class is exhibited below.

\paragraph{The initial coefficient calculation.}
Put $S=\Gamma_{\mathbb Z}(T)$.  Choose $h=2^a$, and set
\[
 x=T^{[h]},\quad y=T^{[2h]},\quad
 x^2=2\kappa y,\quad \kappa=\tfrac12\binom{2h}{h}\text{ odd}.
 \tag{136.CH1}
\]
In $\bar S=S/2$, the binary-digit presentation
$\bar S=\bigotimes_{i\ge0}\mathbb F_2[x_i]/(x_i^2)$,
$x_i=T^{[2^i]}$, gives
\[
                 \operatorname{Ann}_{\bar S}(\bar x)=(\bar x)=I.
 \tag{136.CH2}
\]
This follows by splitting the $a$-th factor as
$\bar S'\oplus\bar x\bar S'$; multiplication by $\bar x$
identifies the first summand with the second and kills the
second.  The presentation and oddness of $\kappa$ follow
directly from $\nu_2(n!)=n-s_2(n)$.

Start with an even disk $v,U$ of forms two and three,
both of weight $2h$, with $dv=2U$.  Its class $xU$ is
nonzero of exact order two.  Attach a free pair $t,c$,
of forms one and two and weight $3h$, with
\[
                 dc=xU,\qquad dt=xv-2c.
 \tag{136.CH3}
\]
Call the resulting algebra $E_1$.  Define
\[
 X=xc-\kappa yv,\quad Y=-xt,\quad Z=-\kappa yt.
 \tag{136.CH4}
\]
Then
\[
 dX=0,\quad dY=2X,\quad dZ=xX,\quad
 xY=2Z,\quad xZ=\kappa yY.
 \tag{136.CH5}
\]
The weights of $X,Y$ are $4h$, and that of $Z$ is $5h$.

Here are the exact low groups, with their coefficient
relations:
\[
 H^0(E_1)=S,\quad H^1(E_1)=0,\quad
 H^2(E_1)\simeq I\langle3h\rangle,\quad
 H^3(E_1)\simeq(S/(2,x))[U].
 \tag{136.CH6}
\]
Indeed a closed two-form is $\alpha v+\beta c$ with
$2\alpha+x\beta=0$.  Its class maps to $\bar\beta\in I$.
Writing $\beta=xr+2s$ shows that this cycle is
$rX-s\,dt$.  Its class is zero precisely when
$r\in(2,x)$, by (136.CH5) or by inspecting the coefficient
of $c$ modulo two.  In form three the basis is
$U,tv,tc$.  The differential of $tc$ has the unique term
$-4\gamma_2(c)$, and after its coefficient vanishes the
differential of $tv$ has the nonzero independent $-2vc$
term.  Thus the three-cycle module is $SU$, with boundaries
$(2,x)U$.  The one-form calculation uses the independent
$-2c$ term of $dt$.  This proves (136.CH6).

Tensor $E_1$ with a second scaled even disk $b,z$, of
forms two and three and weight $n>0$, with $db=2z$.
Write $D$ for this tensor and $W=4h+n$.  Its low groups are
\[
 H^2(D)=I\langle3h\rangle,
 \qquad H^3(D)=(S/(2,x))[U]\oplus(S/2)[z].
 \tag{136.CH7}
\]
For example, the only additional mixed three-form is $tb$,
whose differential has the independent nonzero $-2cb$ term.
This also verifies (136.CH7) by the same basis calculation.

Consider the four-cycle
\[
                         \theta=Xb-Yz,\qquad d(Yb)=2\theta.
 \tag{136.CH8}
\]
Its exact annihilator in $S$ is $(2,x)$.  To see this without
ignoring other classes of $D$, grade the second disk by its
own arity.  Its arity-one subcomplex is
$[Sb\xrightarrow{2}Sz]$ in forms two and three.
The cone exact sequence gives
\[
 0\longrightarrow (H^1(E_1)/2)\langle n\rangle
 \longrightarrow H^4(E_1\otimes[Sb\to Sz])
 \longrightarrow H^2(E_1)[2]\langle n\rangle\longrightarrow0.
 \tag{136.CH9}
\]
The first term is zero; the last is $I\langle3h+n\rangle$.
The class of $\theta$ maps to $[X]$ with this weight shift.
The arity decomposition
is a direct sum of complexes, so no other arity changes its
annihilator.  Equations (136.CH2), (136.CH6) prove the claim.
In particular this is a nonzero class of minimal integer
order two.

\paragraph{The minimal degree-four stage and all old pairs.}
Attach $a,k$, of forms three and two and weight $W$, by
\[
                       da=\theta,\qquad dk=2a-Yb.
 \tag{136.CH10}
\]
Write $D'$ for the result.  Its second cohomology is still
$I\langle3h\rangle$.  The degree-four low-cohomology exact
sequence identifies the quotient of $H^3(D')$ by $H^3(D)$
with the kernel of
$(S/2)\langle W\rangle\xrightarrow{[\theta]}H^4(D)$,
which is $I\langle W\rangle$ by (136.CH9).
An explicit lift of its generator is
\[
 H=xa-Zb,\qquad dH=0,\qquad
                    2H=d(xk),\quad xH=d(\kappa yk).
 \tag{136.CH11}
\]
These identities follow from (136.CH5).  Consequently every
third class has a representative
\[
                        \lambda U+\mu z+\nu H,
 \tag{136.CH12}
\]
and every such class is killed by two.  We need no splitting
assertion: the exact sequence, the generation of $I$ by
$\bar x$, and the displayed lift suffice.  Every second
class is represented by a coefficient multiple of $X$.

All these algebras are rationally vertically exact.  For
$E_1$, replace $c$ rationally by $c-xv/2$, so that
(136.CH3) is a unit disk after rescaling.  The second
scaled disk is rationally a unit disk.  In (136.CH10)
replace $a$ by $a-Yb/2$ to isolate one more such disk.
All coordinate changes are triangular, weight-preserving
free special-PD changes over $\mathbb Q$.

Use the supplied first-skeleton replacement
$\rho:C\to B$, with $C_0=S$, $C_1=F_1$, and its free
envelope $L(F)$.  Define an actual two-truncated source
by $A_0=S$, $A_1=F_1$, and
$A_2=L(F)_2\otimes_S D'$.  Every new generator has zero
face; degeneracies are those of $L(F)$.  The differential
equations have zero faces and thus define all simplicial
identities on this truncation.  The $S$-linear contraction
of $L(F)_2$ carries all preceding cohomology calculations
to $A_2$.

Choose $h$ and $n$ above the finite required guards, with
the $a$-th binary digit of $n$ equal to one.  The same
digit of $W=4h+n$ is one.  Choose a closed relative form
\[
 \eta\in C^{W,2}(\Delta[2],G),\quad
 G=\partial\Delta[2],\qquad
                         \int\rho\eta=T^{[W]}.
 \tag{136.CH13}
\]
The period of $x\eta$ is even, by the binary binomial
calculation.  The relative coefficient comparison in
weight $W+h$ therefore supplies a closed relative two-form
$\chi$ and a relative one-form $\psi$ with
\[
                              d\psi=2\chi-x\eta.
 \tag{136.CH14}
\]
Define two maps, canonical on $L(F)$, zero on every new
generator except $f_1(k)=\eta$; set $f_0(k)=0$.
All relations and faces are respected.

Here is the full old all-modulus verification.  Suppose
$d\xi=0$, $de=N\xi$, with $|\xi|=3$, $|e|=2$.
If $N$ is odd, (136.CH12) implies $\xi=d\alpha$;
then $e-N\alpha=\ell X+d\beta$.  The choices
$c_j=f_j\alpha$, $t_j=-f_j\beta$ give
$dc_j=f_j\xi$, $dt_j=Nc_j-f_je$ and common graph values.
If $N=2m$, write
\[
 \begin{aligned}
 \xi&=\lambda U+\mu z+\nu H+d\alpha,\\
 e&=m(\lambda v+\mu b+\nu xk)+N\alpha+\ell X+d\beta.
 \end{aligned}
 \tag{136.CH15}
\]
The second identity follows because the difference is a
closed two-form and $H^2$ is generated by $[X]$.
With $\varepsilon_0=0$, $\varepsilon_1=1$, take
\[
 c_j=f_j\alpha+\varepsilon_j\nu\chi,
 \qquad t_j=m\varepsilon_j\nu\psi-f_j\beta.
 \tag{136.CH16}
\]
Equations (136.CH14)--(136.CH15) check both differential
identities literally.  Their graph restrictions agree,
since $\chi,\psi$ are relative and the old maps agree on
the supplied first skeleton.  This proves the claimed
property for every old pair, weight and modulus.

\paragraph{A minimal degree-two repair and retained RH1.}
The old class $[X]$ is nonzero of exact integer order two.
Adjoin a full-boundary pair $r,q$ of forms one and zero
and weight $4h$, with
\[
 dr=X,\qquad dq=2r-Y,\qquad r|_G=q|_G=0.
 \tag{136.CH17}
\]
Both maps extend by zero on $r,q$.  Put
\[
             A_*=a-rb-qz,\qquad K_*=k-qb.
 \tag{136.CH18}
\]
Direct differentiation gives $dA_*=0$ and $dK_*=2A_*$.
For \emph{every} retained target extension,
\[
                  f_jA_*=0,\qquad f_jK_*=\varepsilon_j\eta,
 \tag{136.CH19}
\]
because the old $a,b,z$ have zero images.  Common-graph
primitives would give, after applying $\rho$ and Stokes,
\[
                        0=2\int\rho(C_1-C_0)-T^{[W]},
 \tag{136.CH20}
\]
where the integral is in $S^W$.  This is impossible.
The residue is $T^{[W]}/2$ modulo $S$, up to the one
fixed orientation sign.  Multiplying the pair by
$T^{[2^b]}$, with $2^b>W$, preserves nonzero residue:
$\binom{W+2^b}{W}$ is odd.  This gives unbounded output
weights with the same targets.

For precision, low repair is not the identity here.
The one-form $xr-Z$ is closed and
$2(xr-Z)=d(xq)$.  It is not a boundary: after the vertex
augmentation of $L(F)_2$, a zero-form of weight $5h$ in
the free extension is a coefficient plus $s q$ with
$s\in S^h$.  Its $r$-coefficient is $2s$, which cannot
equal $x$.  Higher divided powers of $q$ have weight at
least $8h$.  Thus there is a genuine new one-class.

Nevertheless the hypotheses of AQ-136-RELATIVE-H1 hold.
The simplex zero and structural weight zero are unchanged.
Rationally take the new coordinate $v_0=2r-Y$, so that
$dq=v_0$ and $dv_0=0$, to identify (136.CH17)
with a free even zero-form unit disk over the previous
rational algebra.  Hence rational coefficient comparison
persists uniformly on the finite truncation.  The same RH1
construction uses normalized degrees zero, one and two only
on this truncation, and retains both maps.  It
does not change the already exact first skeleton.  The
images (136.CH19) and the equations for the pair persist
at every step and in the colimit.  Thus the same Stokes
contradiction (136.CH20) applies after RH1, independently
of what other higher cohomology the repair creates or kills.
Choosing the original generator weights above its fixed
guard also makes the untouched lower weights exact through
one, so the repaired truncation has $H^0=S$, $H^1=0$ in
all weights.  No graded-freeness after RH1 is claimed.
\hfill$\square$

\paragraph{Exact history limitation.}
The two initial scaled even disks were specified as input.
The subsequent classes $xU$, $\theta$ and $X$ all have
minimal integer order two, as proved above.  This does
not realize the two initial disks through a minimal
construction starting at $L(F)$, nor force these retained
maps in an arbitrary successful common-source construction.
It refutes preservation from the stated all-pair graph
condition alone even with this limited minimal-order
provenance.

\paragraph{Verification status.}
The author completed a whole saved-proof reading, checking
all low bases, the shifted cone exact sequence and annihilator,
the literal all-modulus primitive formulas, the supplied first
skeleton, both minimal-order claims, the new one-class, and
the retained-map RH1 obstruction.  The missing weight shifts
in the cone display were corrected during this reading.  No
mathematical gap was found.  Reviewer cote123 completed a whole
reading of the current saved proof, including the corrected shifts
and truncated RH1 paragraph, and reported PASS.  Reviewer cote15
also completed a whole saved-proof reading and reported PASS,
including the binary annihilator, all-modulus formulas, literal
first skeleton, new one-class and retained RH1 scope.  Root
completed a whole saved-proof reading, including the corrected
cone grading, and reported PASS.  Its event is recorded centrally
as a contributor whole-proof reading.  The parent proposed this
coefficient-history test; reviewers123 and1348 checked its
communicated annihilator calculation before this complete
proof was written.  Those communications are not counted
as saved-file readings.
